\documentclass[11pt,a4paper]{article}
\usepackage[margin=1in]{geometry}
\usepackage{amsmath,amssymb}
\usepackage{bm}
\usepackage{enumitem}
\usepackage{booktabs}
\usepackage{setspace}
\onehalfspacing

\title{Appendix Audit: Issues and Revision Guidance}
\author{Technical Note}
\date{3 April 2026}

\begin{document}
\maketitle

This note documents issues found during a systematic audit of all four
appendices, combining mechanical findings with substantive exposition
feedback (the latter drawing on a parallel review by Codex). Mechanical
issues have been resolved; substantive revision suggestions remain for a
focused prose pass.

%==============================================================================
\section{Appendix A: Hyperprior Derivation by Distributional Family}
%==============================================================================

\texttt{appendix\_hyperprior\_families.tex} \hfill Label:
\texttt{app:hyperprior\_families}

\subsection{Status: Clean}

No mechanical issues. All cross-references resolve, all 12 figures exist
on disk, style conforms to guidelines.

\subsection{Substantive suggestions}

\paragraph{Gamma-IG figure reuse.}
The Gamma-IG subsection reuses the NIG sensitivity heatmaps for
within/between-population standard deviations rather than producing its
own. This is reasonable since the analytical sensitivity is identical
across NIG and Gamma-IG (the heatmaps depend only on $\rho$, $\kappa$,
and $v^*$, not on the distributional family). However, the text does not
explain this. One sentence noting that the sensitivity surfaces are
family-invariant and therefore shared with the NIG case would make the
reuse feel intentional rather than accidental.

\paragraph{Beta-Beta cross-moment intuition.}
The derivation of $\operatorname{Var}[\mu_k(1-\mu_k)\psi_k]$ becomes algebraically
dense. An extra sentence before this calculation explaining \emph{why}
the cross-moment term arises (the Beta variance is a nonlinear function
of both $\mu$ and $\psi$, so even with independence between the
hyperpriors, the unconditional variance of $z_{ik}$ involves their
product) would help the reader understand that the complication is
structural.

\paragraph{Summary subsection.}
The summary currently emphasizes ordering by KS distance. It could also
state explicitly that the family choice governs two distinct things: the
feasibility of $(\rho, \kappa)$ pairs (each family has a different
feasibility region) and how closely the compound predictive distribution
matches the direct prior away from its first two moments. These are
related but not identical, and separating them would sharpen the
comparative discussion.


%==============================================================================
\section{Appendix B: Injectivity Analysis}
%==============================================================================

\texttt{appendix\_global\_injectivity.tex} \hfill Label:
\texttt{app:global\_injectivity}

\subsection{Mechanical issues (resolved)}

\begin{enumerate}[label=\arabic*.]
  \item Three em-dash violations removed (lines 171, 197, 282--283).
  \item All \texttt{\textbackslash cref\{\}} calls replaced with
    explicit \texttt{\textbackslash ref\{\}} plus type prefix.
  \item All 16 internal labels renamed from \texttt{sec:} to
    \texttt{appB:} namespace.
  \item Dangling Jacobian factorization note (lines 971--975) removed.
\end{enumerate}

\subsection{Substantive suggestions}

This is the most technically demanding appendix (975 lines, the longest
in the paper). The mathematics is correct and the sequential
decomposition strategy is elegant, but the reader load rises sharply in
the middle third. The following additions would improve navigability
without changing any results.

\paragraph{Roadmap before the CRRA proof sequence.}
The appendix establishes three conditions (A, B, C) and then verifies
each for the CRRA-Prelec case. Before the verification begins (around
the start of what is now the Condition~A verification section),
a short paragraph previewing the proof strategy would help. Something
like: ``The remainder of this appendix verifies Conditions A, B, and C
for the CRRA-Prelec specification used in the main text. Condition~A is
established via a two-point injectivity lemma for the Prelec family
(Section~X), which reduces the PWF identification to a diagonal system.
Condition~B is the most involved step: it requires showing that the CRRA
CE function, viewed as a function of $r$ alone after fixing decision
weights, is strictly monotone. This is accomplished through a derivative
identity (Lemma~Y) and a kernel monotonicity argument (Lemma~Z).
Condition~C follows immediately from the Fechner noise structure.''

\paragraph{Interpretive sentence after the derivative identity.}
After the derivative identity lemma (currently Lemma~B.5), the proof has
reduced the Condition~B sign question to comparing two outcome ratios
through a monotone kernel. One sentence stating this strategic
consequence in plain language would prevent the reader from losing the
thread. For example: ``This identity reduces the question of whether
$\text{CE}_1(r) = \text{CE}_1(r')$ to whether a weighted integral of a
kernel function $H_w$ can vanish, which the next lemma rules out by
showing $H_w$ is strictly monotone.''

\paragraph{Interpretive sentence after kernel monotonicity.}
The kernel monotonicity proof (Lemma~B.6) is careful but ends with the
mathematical conclusion $\partial H_w / \partial \rho < 0$. One sentence
explaining why this is the exact property needed for Condition~B
(``Because $H_w$ is strictly decreasing in $\rho$ and the integration
kernel is positive, the integral cannot vanish unless $s = 0$, which
gives $r = r'$'') would connect the lemma back to the proposition it
serves.

\paragraph{Notation slip around $\Delta u$.}
Around lines 107--113 of the original, the wording shifts between a
binary choice task comparing $A$ and $B$ and language referring to
``lottery $k$''. This is minor but tightening it would improve local
clarity.


%==============================================================================
\section{Appendix C: Identification from Choice Probabilities under EUT}
%==============================================================================

\texttt{appendix\_eut\_choice\_probs.tex} \hfill Label:
\texttt{app:eut\_choice\_probs}

\subsection{Mechanical issues (resolved)}

\begin{enumerate}[label=\arabic*.]
  \item Broken \texttt{\textbackslash ref\{sec:framework\}} changed to
    \texttt{\textbackslash ref\{sec:gq\_space\}}.
  \item All 52 occurrences of
    \texttt{\textbackslash boldsymbol\{\textbackslash theta\}} replaced
    with the \texttt{\textbackslash bth} macro.
  \item Three elaborative colons restructured into separate sentences or
    conjunctions.
\end{enumerate}

\subsection{Substantive suggestions}

\paragraph{Positioning relative to Appendix B.}
The opening paragraph now connects to the main text's discussion of
injectivity and scopes the result to EUT. However, the relationship
between this appendix and Appendix~B could be made more explicit.
Appendix~B proves injectivity for a specific RDU model using CEs;
Appendix~C proves injectivity for a broad class of EUT models using
choice probabilities. These are complementary results, not competing
ones. One sentence noting that Appendix~C trades the richer behavioral
functional (RDU) for a cleaner design-side structure (the logistic link
is a diffeomorphism, so the contrast matrix $C$ carries all the
design-side identification content) would help the reader understand the
division of labor.

\paragraph{Part 4 interpretation.}
Part~4 of the theorem (the $m$-to-1 pullback) is one of the more
interesting results because it directly justifies the Bayesian
change-of-variables prior even when the forward map is not globally
injective. An extra sentence of interpretation connecting this to the
paper's framework (``This means the researcher's quantity-space prior is
recovered on the feasible image regardless of whether the forward map
has finitely many preimages, provided the multiplicity is constant'')
would strengthen the link.

\paragraph{``Easy to verify'' claim.}
The opening paragraph promises that the rank condition is ``easy to
verify.'' The appendix gives enough structure for this to be plausible,
but does not quite cash it out. Consider adding a one-sentence remark
after the binary design example noting that for any specific design, the
researcher can check the rank condition numerically by evaluating $C
\cdot D\phi(\bm{\theta}_0)$ at a single convenient parameter value.


%==============================================================================
\section{Appendix D: Sensitivity to $(\rho, \kappa)$}
%==============================================================================

\texttt{appendix\_rk\_sensitivity.tex} \hfill Label:
\texttt{app:rk\_sensitivity}

\subsection{Mechanical issues (resolved)}

\begin{enumerate}[label=\arabic*.]
  \item Subject references reworded from
    ``Subject~\textbackslash ref\{fig:...\}'' to proper figure references.
  \item Stale median updated from 1.09 to 1.08.
  \item Closing synthesis paragraph added (by author).
\end{enumerate}

\subsection{Substantive suggestions}

\paragraph{Pythontex integration.}
All numerical values in this appendix are hardcoded in the \LaTeX{}
source. This makes them fragile against re-runs. The main text avoids
this by loading values from \texttt{paper\_data.npz} via pythontex. The
same approach should be extended here: save the prior predictive medians
and subject classification counts into the NPZ during the pipeline, and
load them with \texttt{\textbackslash py\{D['...']\}}.

\paragraph{Summary table.}
Codex suggested adding a small table with one or two quantitative
metrics per configuration (e.g., prior-dominated share, median absolute
CE shift). The data for this already exists in
\texttt{rk\_sensitivity\_summary.npz}. A 4-row table showing the
classification counts across all four configurations would make the
appendix feel more analytical and give the closing synthesis paragraph
concrete numbers to reference. This could be generated via pythontex
from the summary NPZ.

\paragraph{CDF difference and CE overlay subsections.}
These subsections currently present figures with minimal interpretive
text. Even one sentence per configuration noting the main difference
from the baseline (e.g., ``With $\rho = 0.8$, the prior-dominated share
increases to X\%, reflecting the larger fraction of variance attributed
to mean uncertainty'') would make the appendix read less like a figure
repository.


%==============================================================================
\section{Cross-Cutting Notes}
%==============================================================================

\paragraph{Genre consistency.}
The four appendices span different styles: pedagogical derivation (A),
theorem-proof (B), compact methodological extension (C), and visual
evidence (D). This is not inherently a problem, but slightly stronger
framing at the start of each appendix (one sentence on what the appendix
does and why it is here) would make the package feel more cohesive.

\paragraph{Uninformative model re-run.}
The diffuse NIG model (\texttt{rdu\_bhm\_uninf}) has not been re-run
with float64. Its MCMC data dates from 23 March 2026 (float32). While
the NIG model does not use the CRRA utility function directly in its
posterior (it operates in parameter space with NIG hyperpriors, not
through a CE forward map), the float32 data may still contain minor
numerical artifacts. Before final submission, this model should be
re-run with \texttt{use\_double: True} via \texttt{do\_sampling\_uninf}
in \texttt{main.py}.


%==============================================================================
\section{Summary of Remaining Work}
%==============================================================================

\begin{center}
\begin{tabular}{llll}
  \toprule
  Appendix & Item & Type & Effort \\
  \midrule
  A & Gamma-IG figure reuse note & 1 sentence & Small \\
  A & Beta-Beta cross-moment intuition & 1 sentence & Small \\
  A & Summary: family choice governs two things & 2 sentences & Small \\
  B & Roadmap paragraph before CRRA proofs & 1 paragraph & Medium \\
  B & Interpretive sentence after derivative identity & 1 sentence & Small \\
  B & Interpretive sentence after kernel monotonicity & 1 sentence & Small \\
  B & $\Delta u$ notation slip & Wording & Small \\
  C & Position relative to Appendix B & 1 sentence & Small \\
  C & Part 4 interpretation & 1 sentence & Small \\
  C & ``Easy to verify'' remark & 1 sentence & Small \\
  D & Pythontex integration & Pipeline + \LaTeX & Medium \\
  D & Summary table & Pythontex table & Medium \\
  D & Per-configuration interpretive sentences & 3 sentences & Small \\
  -- & Re-run uninf model with float64 & MCMC & Large \\
  \bottomrule
\end{tabular}
\end{center}

Most items are single-sentence additions that can be done in one focused
pass through the appendices. The two medium items for Appendix~D
(pythontex wiring and summary table) require pipeline changes. The
uninformative model re-run is the largest item but is a single
\texttt{do\_sampling\_uninf = True} run.

\end{document}
